\mainmatter

\chapter{4x10\textsuperscript{9}D.C.}\label{4x10^9}

\hfill\break

Tutti cadiamo, e finiamo a terra da qualche parte.

Così affittammo una camera a Padang, al terzo piano di un hotel in stile
coloniale dove per un po' saremmo passati inosservati.

Novecento euro a notte ci garantirono riservatezza e un balcone con
vista sull'oceano Indiano.

Con il bel tempo, che negli ultimi giorni non era mancato, riuscivamo a
vedere la parte più vicina dell'Arco: una linea verticale del colore
delle nuvole che si levava dall'orizzonte e svaniva, salendo, nella
foschia azzurra. Per quanto fosse impressionante, soltanto una parte
dell'intera struttura si scorgeva dalla costa occidentale di Sumatra.
L'estremità inferiore dell'Arco scendeva verso i picchi sommersi del
Carpenter Ridge, a più di mille chilometri di distanza, attraversando la
Mentawai Trench come una fede nuziale caduta dritta in uno stagno poco
profondo. Sulla terraferma sarebbe andata da Bombay, sulla costa
orientale dell'India, a Madras, verso ovest. Oppure, grossomodo, da New
York a Chicago.

Diane aveva trascorso quasi tutto il pomeriggio sul balcone, sudando
sotto un ombrellone a righe sbiadito. La vista la affascinava, e io ero
contento e sollevato che, dopo tutto ciò che era accaduto, fosse ancora
capace di provare un tale piacere.

La raggiunsi al tramonto. Il tramonto era il momento migliore.

Un mercantile che navigava lungo la costa verso il porto di Teluk Bayur
divenne una collana di luci che scivolava morbida al largo
nell'oscurità. La parte più vicina dell'Arco luccicava come un chiodo
rosso brunito che appuntava il cielo al mare. Guardammo l'ombra della
Terra salire lungo il pilastro mentre la città si oscurava.

Era una tecnologia ``\emph{indistinguibile dalla magia}'', come recita
la famosa citazione. Cos'altro se non la magia avrebbe potuto creare un
flusso d'aria e di mare continuo dal golfo del Bengala all'oceano
Indiano in grado di trasportare una nave verso porti ben più
sconosciuti? Quale miracolo dell'ingegneria consentiva a una struttura
con un raggio di mille chilometri di sostenere il proprio peso? Di cosa
era composta, e come poteva fare ciò che faceva?

Forse solo Jason Lawton avrebbe potuto rispondere a queste domande, ma
Jason non era con noi.

Diane sedeva scomposta su una sdraio, con il prendisole giallo e il
buffo e ampio cappello di paglia ridotti a geometrie di ombre
dall'oscurità che si stava addensando. La sua pelle era lucida, liscia,
color nocciola. Gli occhi catturavano con grande fascino gli ultimi
bagliori, ma la sua espressione era ancora diffidente - quella non era
cambiata.

Sollevò lo sguardo verso di me. «È tutto il giorno che ti agiti.»

«Sto pensando di scrivere qualcosa prima che inizi. Una specie di
memoriale.»

«Hai paura di ciò che potresti perdere, Tyler? Ma è irrazionale. La
memoria non viene cancellata.»

No, cancellata no; ma c'era la possibilità che rimanesse offuscata,
distorta, sfocata. Gli altri effetti collaterali del farmaco erano
momentanei e tollerabili, ma l'eventualità di perdere i ricordi mi
terrorizzava.

«Comunque,» proseguì, «le probabilità giocano a tuo favore. Lo sai
meglio di chiunque altro. \emph{C'è} il rischio\ldots{} ma è \emph{solo}
un rischio, e anche piuttosto basso, direi.»

E se fosse capitato a lei, forse sarebbe stata una benedizione.

«In ogni caso,» replicai, «mi sentirei meglio se scrivessi qualcosa.»

«Se non vuoi proseguire, non devi farlo per forza. Quando sarai pronto
lo capirai.»

«No, voglio farlo.» O almeno così dissi a me stesso.

«Allora dobbiamo iniziare stasera.»

«Lo so, ma nelle prossime settimane\ldots»

«Probabilmente non avrai voglia di scrivere.»

«A meno che non ne possa fare a meno.» La grafomania era un altro dei
possibili effetti collaterali, ma il meno angosciante.

«Aspetta a dirlo quando ti verrà la nausea.» Mi rivolse un sorriso
consolatorio. «Tutti abbiamo paura di doverci lasciare alle spalle
qualcosa.»

Era vero, ma la cosa mi metteva ansia e non ci volevo pensare.

«Senti,» dissi, «forse dovremmo davvero iniziare.»

L'aria odorava di tropici, ma con una venatura di cloro proveniente
dalla piscina dell'hotel, tre piani più in basso. A quel tempo Padang
era un grande porto internazionale, zeppo di stranieri: indiani,
filippini, coreani e anche americani raminghi come me e Diane, gente che
non poteva permettersi un passaggio di lusso e non era qualificata per i
programmi di reinsediamento approvati dall'ONU. Era una città vivace ma
spesso illegale, specialmente da quando a Jakarta era salita al potere
la nuova Reformasi.

L'hotel era comunque sicuro e le stelle erano là fuori in tutta la loro
gloria punteggiata. La sommità dell'Arco in quel momento era la cosa più
luminosa del cielo, una delicata U d'argento scritta sottosopra da un
dio dislessico. Mentre la guardavamo svanire, stringevo la mano di
Diane.

«A cosa pensi?» mi chiese.

«All'ultima volta che ho visto le vecchie costellazioni.» Vergine,
Leone, Sagittario: il lessico dell'astrologo ridotto a note a piè di
pagina di un libro di storia.

«Ma da qui, nell'emisfero meridionale, sarebbero state diverse, vero?»

Supponevo di sì.

Poi, nell'oscurità totale della notte, tornammo nella stanza. Accesi le
luci mentre Diane tirava le tende e spacchettava la siringa e il kit di
fiale che le avevo insegnato a usare. Riempì la siringa sterile e,
concentrandosi, la picchiettò per far uscire l'aria. Sembrava esperta ma
le tremava la mano.

Mi tolsi la camicia e mi distesi sul letto.

«Tyler\ldots» all'improvviso era lei a essere riluttante.

«Nessun ripensamento» dissi. «So a cosa vado incontro. E ne abbiamo
parlato una decina di volte.»

Annuì e mi strofinò l'interno del gomito con l'alcool. Teneva la siringa
nella mano destra, con la punta verso l'alto. La piccola quantità di
fluido all'interno sembrava inoffensiva come l'acqua.

«Fu molto tempo fa» mormorò.

«Cosa?»

«Quella volta che guardammo le stelle.»

«Sono contento che tu non l'abbia dimenticato.»

«Certo che non l'ho dimenticato. Ora stringi la mano a pugno.»

Il dolore fu leggero. Per lo meno all'inizio.